\documentclass{article}
\usepackage{amsmath, amssymb, booktabs}

\title{Frozen-Jacobian Line Search for the Geodesic IPM}
\author{}
\date{}

\begin{document}
\maketitle

\section{Setup}

We consider the geodesic interior-point method for the conic program
\[
  \min_x\; c^T x \quad \text{s.t.}\quad Ax + b \succeq 0,
\]
with the weight parameterization
\[
  s = \tfrac{1}{k}\,W^{-1}(e - d), \qquad
  \lambda = \tfrac{1}{k}\,W(e + d),
\]
where $k = 1/\sqrt{\mu}$ and $W$ is a positive diagonal scaling (or,
more generally, an element of the automorphism group of the cone).
We write $P(W) = W^2$ for the quadratic representation.

Eliminating $d$ from the KKT conditions $s = Ax + b$ and
$A^T\lambda = c$ yields the normal equations
\begin{equation}\label{eq:standard}
  (A^T W^2 A + Q)\,y = 2A^T W - k(c + A^T W^2 b),
\end{equation}
where $y = kx$ and $Q$ accounts for any quadratic cost.  The Newton
direction is recovered as $d = e + W\,\text{slack}$, where
$\text{slack} = -kb - Ay$.

The standard geodesic LP algorithm repeats:
\begin{enumerate}
\item Factor $G = A^T W^2 A + Q$.
\item Decompose into $d = d_0 + k\,d_1$ (2 back-solves).
\item Line search: find $k_{\max}$ with $\|d_0 + k_{\max} d_1\|_\infty \le 1$.
\item Geodesic step: $W \leftarrow W \cdot \exp(\alpha\, d)$.
\end{enumerate}
Each outer iteration requires one factorization and two back-solves.

\section{Frozen-Jacobian Newton system}

When we step from $W_0$ (the factorization point) to a new iterate $W_i$,
the standard approach refactors at $W_i$.  Instead, we freeze the Jacobian
at $W_0$ and only update the residual at $W_i$.

The standard linearized equations at $W$ with direction $d$ are:
\begin{align}
  A^T\bigl(W(e+d)\bigr)/k &= Qx + c, \label{eq:dual-std}\\
  W^{-1}(e-d)/k &= Ax + b. \label{eq:primal-std}
\end{align}
In the frozen-Jacobian variant, the ``value'' terms use $W_i$ (current
iterate) while the ``Jacobian'' terms use $W_0$ (frozen):
\begin{align}
  A^T(W_i + W_0\, d)/k &= Qx + c, \label{eq:dual-frozen}\\
  (W_i^{-1} - W_0^{-1}\, d)/k &= Ax + b. \label{eq:primal-frozen}
\end{align}

\subsection{Elimination}

From~\eqref{eq:primal-frozen}:
\[
  d = W_0\bigl(W_i^{-1} - k(Ax + b)\bigr)
    = W_0 W_i^{-1} - k\,W_0(Ax + b).
\]
Substituting into~\eqref{eq:dual-frozen}:
\[
  A^T\!\bigl(W_i + W_0^2 W_i^{-1} - k\,W_0^2(Ax+b)\bigr)/k = Qx + c.
\]
Rearranging:
\begin{equation}\label{eq:frozen-system}
  \boxed{
  (A^T W_0^2 A + Q)\, x
  = \frac{1}{k}\,A^T\!\bigl(W_i + W_0^2 W_i^{-1}\bigr)
    - A^T W_0^2 b - c.
  }
\end{equation}

\paragraph{Consistency check.}
When $W_i = W_0$, the centering term becomes
$W_0 + W_0^2 W_0^{-1} = 2W_0$, recovering~\eqref{eq:standard}.

\subsection{Direction recovery}

After solving~\eqref{eq:frozen-system} for $y = kx$:
\begin{equation}\label{eq:d-recovery}
  d = \frac{W_0}{W_i} + W_0 \cdot \text{slack},
  \qquad \text{slack} = -kb - Ay.
\end{equation}
Compare with the standard recovery $d = e + W \cdot \text{slack}$:
the centering term changes from $e$ to $W_0/W_i$, and the Jacobian
in the slack term changes from $W_i$ to $W_0$.

\section{Two-term decomposition}

The frozen system~\eqref{eq:frozen-system} decomposes into centering and
cost components:
\begin{align}
  (A^T W_0^2 A + Q)\,y_0 &= A^T(W_i + W_0^2 W_i^{-1}),
    &\quad &\text{(centering, depends on $W_i$)} \label{eq:rhs0}\\
  (A^T W_0^2 A + Q)\,y_1 &= -(c + A^T W_0^2 b) + d_{\text{eq}},
    &\quad &\text{(cost, frozen)} \label{eq:rhs1}
\end{align}
where $d_{\text{eq}}$ is the equality constraint correction.

\paragraph{Key observation.}
The cost component $y_1$ depends only on $W_0$ and is \emph{constant}
across all frozen-Jacobian iterations.  After the initial 2-solve
decomposition, each subsequent iteration requires only \textbf{one
back-solve} (for $y_0$) plus a line search on
$d = d_0^{(i)} + k\,d_1^{\text{frozen}}$.

The directions are recovered componentwise:
\begin{align}
  d_0 &= W_0 \cdot (W_i^{-1} - A y_0), \\
  d_1 &= W_0 \cdot (-b - A y_1). \quad \text{(frozen)}
\end{align}

\section{Algorithm}

\begin{enumerate}
\item \textbf{Factor:} Set $W_0 \leftarrow W_i$.  Compute $G = A^T W_0^2 A + Q$
  and factor.  Decompose $d = d_0 + k\,d_1$ via standard 2-solve decomposition.
  Line search for $k$.  Take geodesic step.
\item \textbf{Frozen-Jacobian iterations} (repeat until degradation):
  \begin{enumerate}
  \item Compute $d_0^{(i)}$ via 1 back-solve with RHS from~\eqref{eq:rhs0}.
    Reuse $d_1^{\text{frozen}}$ from step~1.
  \item Line search: find $k_{\max}$ with
    $\|d_0^{(i)} + k_{\max}\,d_1^{\text{frozen}}\|_\infty \le 1$.
    Set $k \leftarrow \max(k, k_{\max})$.
  \item Take geodesic step $W_i \leftarrow W_i \cdot \exp(\alpha\, d)$.
  \item If $\|d_0^{(i)}\|$ is growing (frozen Jacobian degrading), exit.
  \end{enumerate}
\item Go to step~1 (refactor at new $W_i$).
\end{enumerate}

\paragraph{Cost per outer iteration.}
One factorization, $2 + n_{\text{inner}}$ back-solves (2 for the initial
decomposition, 1 per frozen-Jacobian step), and $1 + n_{\text{inner}}$
line searches.

\section{Results}

On a test QP (HS21-type, 7 variables, 10 constraints):

\begin{center}
\begin{tabular}{lrrr}
\toprule
Method & Factorizations & Back-solves & Final gap \\
\midrule
Standard GeodesicLP & 16 & 32 & $9.7 \times 10^{-11}$ \\
Frozen-Jacobian ($n_{\text{inner}}=3$) & 9 & 50 & $8.4 \times 10^{-12}$ \\
\bottomrule
\end{tabular}
\end{center}

The frozen-Jacobian variant halves the number of factorizations at
the cost of additional back-solves.  For problems where factorization
dominates (large sparse systems), this is a significant win.

\subsection{Convergence behavior}

The gap per factorization shows accelerating convergence:

\begin{center}
\begin{tabular}{crrc}
\toprule
Fac & $k$ & Gap & $k$ ratio \\
\midrule
0 & 1.7 & $2.8$ & --- \\
1 & 4.1 & $6.5 \times 10^{-1}$ & 2.4$\times$ \\
2 & 8.2 & $4.2 \times 10^{-2}$ & 2.0$\times$ \\
3 & 30 & $6.0 \times 10^{-3}$ & 3.7$\times$ \\
4 & 136 & $4.0 \times 10^{-4}$ & 4.5$\times$ \\
5 & 577 & $1.2 \times 10^{-5}$ & 4.2$\times$ \\
6 & $7.3\times10^3$ & $3.0 \times 10^{-8}$ & 12.7$\times$ \\
7 & $4.4\times10^4$ & $1.8 \times 10^{-9}$ & 6.0$\times$ \\
8 & $7.5\times10^5$ & $8.4 \times 10^{-12}$ & 17.2$\times$ \\
\bottomrule
\end{tabular}
\end{center}

The $k$-ratio (hence $\mu$-reduction) per factorization accelerates as
the algorithm approaches the solution.  Near the optimum, $d_1$ is small
(the cost direction vanishes), so the frozen Jacobian stays accurate over
many steps and the line search can push $k$ much further before the
approximation degrades.  This produces near-quadratic convergence in
the number of factorizations.

\subsection{Why frozen-Jacobian centering alone fails}

An earlier experiment froze the Jacobian but held $k$ fixed (pure
centering, no line search).  This reduced $\|d\|_\infty$ from $1.0$
to $\sim 0.6$ per step, but the subsequent fresh decomposition showed
the same $\|d_0\| \approx 1$ regardless.  The centering did not help
the next line search because $d_0$ (the centering residual at a fresh
factorization) is determined by the distance from $W$ to the central
path at the \emph{current} $k$---not by how well centered $W$ was at
the \emph{previous} $k$.

The crucial ingredient is combining the frozen-Jacobian with a line
search: this advances $k$ (reduces $\mu$) using only cheap back-solves,
avoiding the factorization cost of the standard approach.

\section{Extension to $\theta$-continuation}

The $\theta$-continuation algorithm (see \texttt{geodesic\_newton\_direction.tex})
uses a three-term decomposition
\[
  d(k,\tau,\theta) = d_0 + k\bigl(\tau\,d_1^{(0)} + \theta\,d_1^{(\theta)}\bigr)
\]
and selects $(\tau, \theta)$ jointly via the duality identity (gap equation)
and a normalization equation.  Here we derive the frozen-Jacobian variant.

\subsection{Frozen-Jacobian primal-dual equations}

At frozen Jacobian $W_0$ and current iterate $W_i$, the $\theta$-blended
primal-dual equations are:
\begin{align}
  \lambda &= \frac{1}{k}\bigl(W_i + P(W_0^{1/2})(d)\bigr),
    \label{eq:lam-frozen-theta} \\
  s &= \frac{1}{k}\bigl(W_i^{-1} - P(W_0^{-1/2})(d)\bigr),
    \label{eq:s-frozen-theta}
\end{align}
with primal feasibility $s = Ax + \tau b + \theta(e - b)$ and dual
stationarity $A^T\lambda = \tau c + \theta(A^T e - c) + d_{\text{eq}}$.

Comparing with the standard parameterization where
$\lambda = \frac{1}{k}P(W^{1/2})(e + d)$:
the \emph{value} $W_i$ replaces $P(W^{1/2})(e)$,
and the \emph{Jacobian} $P(W_0^{1/2})$ replaces $P(W^{1/2})$.

\subsection{Three-term decomposition}

Eliminating $d$ from~\eqref{eq:s-frozen-theta} as in Section~2.1, the
Gram system $(A^T P(W_0) A + Q)\,y = \text{rhs}$ has three right-hand sides:
\begin{align}
  \text{rhs}_0 &= A^T\bigl(W_i + P(W_0)(W_i^{-1})\bigr),
    &\quad &\text{(centering, depends on $W_i$)} \\
  \text{rhs}_1 &= -\bigl(c + A^T P(W_0)\,b\bigr) + d_{\text{eq}},
    &\quad &\text{(cost, frozen)} \\
  \text{rhs}_2 &= -\text{rhs}_1 - A^T\bigl(e + P(W_0)\,e\bigr),
    &\quad &\text{($\theta$ correction, frozen)}
\end{align}
yielding solutions $y_0, y_1^{(0)}, y_1^{(\theta)}$ with $y = y_0/k
+ \tau\,y_1^{(0)} + \theta\,y_1^{(\theta)}$.

The directions are recovered using $P(W_0^{1/2})$:
\begin{align}
  d_0 &= P(W_0^{1/2})\bigl(W_i^{-1} - A\,y_0\bigr),
    \label{eq:d0-frozen-theta} \\
  d_1^{(0)} &= P(W_0^{1/2})\bigl(-b - A\,y_1^{(0)}\bigr),
    \quad\text{(frozen)} \label{eq:d10-frozen-theta} \\
  d_1^{(\theta)} &= P(W_0^{1/2})\bigl(b - e - A\,y_1^{(\theta)}\bigr).
    \quad\text{(frozen)} \label{eq:d1t-frozen-theta}
\end{align}

\paragraph{Key observation.}  Only $\text{rhs}_0$ (and hence $d_0$)
depends on $W_i$.  The cost and theta components $d_1^{(0)},
d_1^{(\theta)}$ are frozen.  After the initial 3-solve decomposition,
each frozen-Jacobian step requires \textbf{one back-solve} to refresh~$d_0$.

\paragraph{Consistency check.}  When $W_i = W_0$:
$P(W_0)(W_0^{-1}) = e$, so $\text{rhs}_0 = A^T(W_0 + e)$\ldots
no, that is wrong.  We have $W_0 + P(W_0)(W_0^{-1}) = W_0 + W_0 = 2W_0$,
recovering the standard $\text{rhs}_0 = A^T(2W)$.  And
$P(W_0^{1/2})(W_0^{-1}) = W_0 \cdot W_0^{-1} = e$, recovering
$d_0 = e - P(W_0^{1/2})(Ay_0)$. \checkmark

\subsection{Duality identity with frozen Jacobian}

The duality identity (gap equation) from \texttt{geodesic\_newton\_direction.tex}
\S6 is
\begin{equation}\label{eq:gap-identity}
  b^T\lambda + c^Tx + d^T\nu + \frac{x^TQx}{\tau} + \frac{\mu}{\tau}
  = \theta\,R,
\end{equation}
where $R = \langle b, e\rangle + 1$.  With the frozen Jacobian, $\lambda$
and $x$ are:
\begin{align}
  \lambda &= \frac{1}{k}\bigl(W_i + P(W_0^{1/2})(d)\bigr)
    = \frac{W_i}{k}
      + \tau\,P(W_0^{1/2})\,d_1^{(0)}
      + \theta\,P(W_0^{1/2})\,d_1^{(\theta)}
      + \frac{1}{k}\,P(W_0^{1/2})(d_0), \\
  x &= \frac{y_0}{k} + \tau\,y_1^{(0)} + \theta\,y_1^{(\theta)}.
\end{align}
Both are affine in $(\tau, \theta)$.  The key difference from the standard
case is in the $b^T\lambda$ term: the constant part is $b^T W_i / k$
(depends on $W_i$) rather than $b^T P(W^{1/2})(e + d_0)/k$ (depends on
a single~$W$).

\paragraph{Decomposing $b^T\lambda$.}
\begin{align}
  b^T\lambda &= \underbrace{\frac{1}{k}\bigl(b^T W_i
    + b^T P(W_0^{1/2})(d_0)\bigr)}_{\sigma_0}
    + \underbrace{b^T P(W_0^{1/2})\,d_1^{(0)}}_{\sigma_1}\,\tau
    + \underbrace{b^T P(W_0^{1/2})\,d_1^{(\theta)}}_{\sigma_\theta}\,\theta.
    \label{eq:sigma-frozen}
\end{align}
Compare with the standard case where
$\sigma_0 = \frac{1}{k} b^T P(W^{1/2})(e + d_0)$:
the frozen version splits the ``value'' $W_i$ from the ``Jacobian''
$P(W_0^{1/2})$.

The $\sigma_1$ and $\sigma_\theta$ terms depend only on $W_0$ (frozen).
Only $\sigma_0$ changes when $W_i$ is updated.

\paragraph{The $c^Tx$ and $d^T\nu$ terms.}
These use $y_0, y_1^{(0)}, y_1^{(\theta)}$ from the frozen Gram solve:
\begin{align}
  c^Tx &= \frac{c^T y_0}{k} + c^T y_1^{(0)}\,\tau + c^T y_1^{(\theta)}\,\theta
         = \frac{\gamma_0^{(c)}}{k} + \gamma_1\,\tau + \gamma_\theta\,\theta, \\
  d^T\nu &\text{ enters through $d_{\text{eq}}$ in the same pattern.}
\end{align}
The $\gamma_1$ and $\gamma_\theta$ terms use the frozen $y_1^{(0)}$
and $y_1^{(\theta)}$.  Only $\gamma_0^{(c)} = c^T y_0$ changes when
$d_0$ is refreshed.

\paragraph{Quadratic cost.}
The $x^TQx/\tau$ term decomposes identically to the standard case
(see \texttt{geodesic\_newton\_direction.tex} \S6.2), since
$x = y_0/k + \tau\,y_1^{(0)} + \theta\,y_1^{(\theta)}$ has the same
structure.  The frozen Jacobian only affects how $y_0$ is computed,
not the algebraic form.

\subsection{Violation quadratic}

Following the same elimination as the standard case
(\texttt{geodesic\_newton\_direction.tex} \S8), we use the normalization
equation to express $\theta(\tau) = a_0 + a_1\,\tau$, substitute into
the gap equation, and obtain
\begin{equation}\label{eq:frozen-tau-quad}
  \beta\,\tau^2 + \gamma\,\tau + \mu_{\text{eff}} = 0,
\end{equation}
where the coefficients have the same algebraic form as the standard case
but with:
\begin{itemize}
  \item $\sigma_0$ computed from $W_i$ and $W_0$ via~\eqref{eq:sigma-frozen},
  \item $\sigma_1, \sigma_\theta$ computed from $W_0$ only (frozen),
  \item $y_0$ from the frozen Gram solve (refreshed), $y_1^{(0)},
    y_1^{(\theta)}$ frozen.
\end{itemize}

\paragraph{What changes per frozen-Jacobian step.}
After refreshing $d_0$ and $y_0$ (1 back-solve):
\begin{enumerate}
  \item Recompute $\sigma_0 = \frac{1}{k}(b^T W_i + b^T P(W_0^{1/2})(d_0))$
    and $\gamma_0^{(c)} = c^T y_0$.
  \item Recompute the quadratic cost terms involving $y_0$:
    $f = y_0/k + \theta\,y_1^{(\theta)}$, $f^TQf$, $f^TQy_1^{(0)}$.
  \item Solve the quadratic~\eqref{eq:frozen-tau-quad} for $\tau$,
    recover $\theta$.
  \item Binary search over $\theta$ (or direct selection) as in the
    standard algorithm.
\end{enumerate}
The normalization coefficients $N_\tau, N_\theta$ also have a $W_i$-dependent
part (through $\lambda_0$) that must be refreshed.

\subsection{Algorithm for frozen-Jacobian $\theta$-continuation}

\begin{enumerate}
\item \textbf{Factor:} Set $W_0 \leftarrow W_i$.  Compute
  $G = A^T P(W_0) A + Q$ and factor.  Decompose into
  $(d_0, d_1^{(0)}, d_1^{(\theta)})$: 3 back-solves.
  Select $(\tau, \theta)$ via gap + normalization.  Take geodesic step.
\item \textbf{Frozen-Jacobian step:}
  \begin{enumerate}
  \item Refresh $d_0$ via 1 back-solve with
    $\text{rhs}_0 = A^T(W_i + P(W_0)(W_i^{-1}))$.
  \item Recompute the $W_i$-dependent duality coefficients ($\sigma_0$,
    normalization constant, quadratic cost terms involving $y_0$).
  \item Select $(\tau, \theta)$ from the updated quadratic.
  \item Take geodesic step.
  \end{enumerate}
\item Go to step~1 (refactor at new $W_i$).
\end{enumerate}

\paragraph{Cost per outer iteration.}
One factorization, $3 + n_{\text{inner}}$ back-solves (3 for the initial
decomposition, 1 per frozen-Jacobian step), and $1 + n_{\text{inner}}$
$(\tau,\theta)$ selections.  The frozen components ($d_1^{(0)},
d_1^{(\theta)}, \sigma_1, \sigma_\theta, \gamma_1, \gamma_\theta$,
normalization coefficients $N_\tau, N_\theta$) are computed once and
reused.

\end{document}
